% ============================================================
%  customised-commands.tex  —  merged from all venue variants
%  Variants: base, ACL, ACM (acmart), LNCS (llncs),
%            Elsevier (elsarticle), ICLR (iclr2024/iclrworkshop),
%            IEEE (IEEEtran), ICML (icml2024), NeurIPS (neurips_2024),
%            AAAI (aaai24/aaai25), USENIX (usenix)
%
%  Auto-detection uses \@ifclassloaded{<class>}{<yes>}{<no>}.
%  All venue-specific logic is resolved at compile time — no
%  manual editing needed when switching document class.
% ============================================================

\makeatletter

% ============================================================
%  VENUE BOOLEANS  (set once, used throughout)
% ============================================================
\newif\if@isACM    \@isACMfalse
\newif\if@isLNCS   \@isLNCSfalse
\newif\if@isELS    \@isELSfalse
\newif\if@isACL    \@isACLfalse
\newif\if@isICLR   \@isICLRfalse
\newif\if@isIEEE   \@isIEEEfalse
\newif\if@isICML   \@isICMLfalse
\newif\if@isNIPS   \@isNIPSfalse
\newif\if@isAAAI   \@isAAAIfalse
\newif\if@isUSENIX \@isUSENIXfalse

\@ifclassloaded{acmart}         {\@isACMtrue}{}
\@ifclassloaded{llncs}          {\@isLNCStrue}{}
\@ifclassloaded{elsarticle}     {\@isELStrue}{}
\@ifclassloaded{acl}            {\@isACLtrue}{}
\@ifclassloaded{acl2023}        {\@isACLtrue}{}
\@ifclassloaded{acl_natbib}     {\@isACLtrue}{}
\@ifclassloaded{iclr2024}       {\@isICLRtrue}{}
\@ifclassloaded{iclrworkshop}   {\@isICLRtrue}{}
\@ifclassloaded{IEEEtran}       {\@isIEEEtrue}{}
\@ifclassloaded{icml2024}       {\@isICMLtrue}{}
\@ifclassloaded{neurips_2024}   {\@isNIPStrue}{}
\@ifclassloaded{aaai24}         {\@isAAAItrue}{}
\@ifclassloaded{aaai25}         {\@isAAAItrue}{}
\@ifclassloaded{usenix}         {\@isUSENIXtrue}{}

\makeatother

% -------  optional / review package  -----------------------
% \usepackage{easyReview}

% -------  language  ----------------------------------------
% IEEE, ICML, NeurIPS handle language internally — skip babel there.
\makeatletter
\if@isIEEE\else\if@isICML\else\if@isNIPS\else
  \usepackage[english]{babel}
\fi\fi\fi
\makeatother

% -------  geometry  ----------------------------------------
% ACL loads geometry explicitly; other classes manage their own margins.
\makeatletter
\if@isACL
  \usepackage{geometry}
  % \usepackage[letterpaper,top=2cm,bottom=2cm,left=3cm,right=3cm,marginparwidth=1.75cm]{geometry}
\fi
\makeatother

% -------  bibliography  ------------------------------------
% ACM, IEEE, ICML, NeurIPS, AAAI all provide their own citation machinery.
% ICLR: natbib loaded by the class itself; only the \cite redefinition is needed.
% USENIX: uses numeric natbib style (no author-year \cite redefinition).
\makeatletter
\if@isACM\else\if@isIEEE\else\if@isICML\else\if@isNIPS\else\if@isAAAI\else
  \if@isUSENIX
    \usepackage[numbers]{natbib}
  \else\if@isICLR
    % natbib provided by iclr class; redefine \cite only
    \let\oldcite\cite
    \renewcommand{\cite}[1]{\citep{#1}}
  \else
    \usepackage{natbib}  % \bibliographystyle{plainnat}
    \let\oldcite\cite
    \renewcommand{\cite}[1]{\citep{#1}}
  \fi\fi
\fi\fi\fi\fi\fi
\makeatother
% \usepackage[backend=biber,style=numeric-comp,sorting=none]{biblatex}

% -------  graphics & layout  -------------------------------
\usepackage{graphicx}
% IEEE, ICML, NeurIPS load microtype themselves or are incompatible.
\makeatletter
\if@isIEEE\else\if@isICML\else\if@isNIPS\else
  \usepackage{microtype}
\fi\fi\fi
\makeatother
% \usepackage{subfigure}
\usepackage{booktabs}
% AAAI discourages float / wrapfig / ulem.
\makeatletter
\if@isAAAI\else
  \usepackage{float}
  \usepackage{wrapfig}
  \usepackage[normalem]{ulem}
\fi
\makeatother

% -------  USENIX-specific layout packages  -----------------
\makeatletter
\if@isUSENIX
  \usepackage{url}
  \usepackage{cuted}    % strip environment for full-width floats
  \usepackage{capt-of}  % captions outside float environments
\fi
\makeatother

% -------  math  --------------------------------------------
\usepackage{bm}
\usepackage{textcomp}
% ACM, IEEE, ICML, NeurIPS already load amssymb via their font/class setup.
\makeatletter
\if@isACM\else\if@isIEEE\else\if@isICML\else\if@isNIPS\else
  \usepackage{amssymb}
\fi\fi\fi\fi
\makeatother
\usepackage{amsmath}
% LNCS provides its own theorem environments including amsthm — skip it there.
\makeatletter
\if@isLNCS\else
  \usepackage{amsthm}
\fi
\makeatother
\usepackage{amsfonts}
\usepackage{mathtools}
\usepackage{pifont}

% -------  algorithms  --------------------------------------
\usepackage{algorithm}
% ACM and AAAI use algpseudocode; all other venues use algorithmic.
\makeatletter
\if@isACM
  \usepackage{algpseudocode}
\else\if@isAAAI
  \usepackage{algpseudocode}
\else
  \usepackage{algorithmic}
\fi\fi
\makeatother

% LNCS-only: inline comment command for algorithm floats.
\makeatletter
\if@isLNCS
  \newcommand{\algcomment}[2][gray]{%
    \unskip\hfill{\footnotesize\textcolor{#1}{// #2}}%
  }
\fi
\makeatother

% -------  fonts & symbols  ---------------------------------
% ACM and AAAI manage their own fonts — \usepackage{times} conflicts.
\makeatletter
\if@isACM\else\if@isAAAI\else
  \usepackage{times}
\fi\fi
\makeatother
% \usepackage{latexsym}
% AAAI omits gensymb and textgreek (class conflicts).
\makeatletter
\if@isAAAI\else
  \usepackage{gensymb}
\fi
\makeatother

% -------  table & figure utilities  ------------------------
\usepackage{multirow}
% AAAI omits lscape.
\makeatletter
\if@isAAAI\else
  \usepackage{lscape}
\fi
\makeatother
\usepackage{subcaption}

% -------  frames, boxes, hyperlinks  -----------------------
\usepackage{mdframed}
% AAAI disallows hyperref in submissions.
\makeatletter
\if@isAAAI\else
  \usepackage{hyperref}
\fi
\makeatother
\makeatletter
\if@isAAAI\else
  \usepackage{textgreek}
\fi
\makeatother

% -------  color & table color  -----------------------------
\usepackage{tabularx, colortbl}
\usepackage[dvipsnames]{xcolor}

% ACM-only: TikZ physics shorthand color.
% xcol (steel blue) -> RoyalBlue (dvipsnames)
\makeatletter
\if@isACM
  \colorlet{xcol}{RoyalBlue}
\fi
\makeatother

% AAAI-only: large diverging color palette mapped to xcolor dvipsnames/mixes.
\makeatletter
\if@isAAAI
  % --- Diverging palette (main 8) ---
  \colorlet{4A7BB7}{SteelBlue!80!white}         % mid steel blue
  \colorlet{91C4DE}{SkyBlue!70!white}            % light sky blue
  \colorlet{FED081}{Goldenrod!70!white}          % soft gold
  \colorlet{F7834D}{BurntOrange!80!white}        % warm orange
  \colorlet{8ECFC9}{TealBlue!50!white}           % soft teal
  \colorlet{FFBE7A}{Goldenrod!50!white}          % pale amber
  \colorlet{FA7F6F}{Salmon!90!white}             % warm salmon
  \colorlet{82B0D2}{SteelBlue!55!white}          % pale steel blue
  % --- Reds & Pinks ---
  \colorlet{D62728}{red!80!black}                % amaranth red
  \colorlet{E13333}{Red}                         % strong crimson
  \colorlet{C63D52}{BrickRed!80!black}           % dark scarlet
  \colorlet{F67088}{Salmon}                      % coral
  \colorlet{F77189}{Salmon!90!Pink}              % ultrared
  \colorlet{FFB3B2}{Salmon!40!white}             % light strawberry
  % --- Magentas & Purples ---
  \colorlet{BF3E87}{Magenta!75!black}            % magenta rose
  \colorlet{F584C5}{Magenta!50!white}            % tearose
  \colorlet{E377C2}{Orchid}                      % middle purple
  \colorlet{9467BD}{MediumOrchid!80!black}       % soft purple (fallback: Orchid!80!blue)
  \colorlet{7B7AB7}{Periwinkle!80!gray}          % ube lavender
  \colorlet{1B2278}{NavyBlue!90!black}           % royal blue-purple
  \colorlet{4E518B}{MidnightBlue!60!Periwinkle}  % twilight indigo
  % --- Blues ---
  \colorlet{005397}{NavyBlue}                    % deep navy cyan
  \colorlet{023047}{MidnightBlue}                % prussian blue
  \colorlet{1C3168}{NavyBlue!80!black}           % catalina blue
  \colorlet{327BB7}{SteelBlue}                   % steel blue
  \colorlet{1F77B4}{RoyalBlue!80!cyan}           % mid blue
  \colorlet{3BA3EC}{Cerulean!80!white}           % dodger blue
  \colorlet{90C9E6}{SkyBlue!60!white}            % cornflower
  \colorlet{94C3E1}{SkyBlue!65!white}            % light sky blue
  \colorlet{A2B1C3}{CadetBlue!50!white}          % steel gray-blue
  % --- Teals & Greens ---
  \colorlet{0192A8}{TealBlue}                    % cerulean
  \colorlet{17BECF}{Cyan!70!TealBlue}            % battery charged blue
  \colorlet{219EBC}{CadetBlue!80!Cyan}           % pacific blue
  \colorlet{36ADA4}{TealBlue!70!white}           % keppel teal
  \colorlet{2CA02C}{ForestGreen}                 % wageningen green
  \colorlet{AEDB89}{OliveGreen!40!white}         % pale green apple
  \colorlet{B3FEAE}{green!30!white}              % minty pastel
  % --- Yellows & Yellow-Greens ---
  \colorlet{BCBD22}{YellowGreen!80!black}        % acid green
  \colorlet{C99E31}{Goldenrod!85!black}          % tussock mustard
  \colorlet{F0CF61}{Goldenrod!70!Yellow}         % golden yellow
  \colorlet{BC9F48}{Goldenrod!60!brown}          % domino gold
  % --- Browns & Earth Tones ---
  \colorlet{8C564B}{RawSienna!80!black}          % spicy mix brown
  \colorlet{BEA1A5}{RosyBrown}                   % dusty mauve
  % --- Pastels & Neutrals ---
  \colorlet{BFB5D7}{Lavender}                    % pale lavender
  \colorlet{FFD8B0}{Peach!80!white}              % peachy pink
  \colorlet{F3EEE3}{Apricot!30!white}            % merino cream
  \colorlet{D9D9D9}{lightgray}                   % light silver
  % --- Blacks ---
  \colorlet{000000}{black}                       % pure black
  \colorlet{333333}{black!80!white}              % near black
\fi
\makeatother

% ACL table header/row shading — available color names for use in tables:
%   headerpink  -> Salmon!30    headergreen -> LimeGreen!30
%   roworange   -> Peach!50     rowgray     -> lightgray!50
%   textorange  -> BurntOrange

% -------  column types & highlight commands  ---------------
% Color mapping from originals to xcolor dvipsnames:
%   p13 (soft lavender) -> Lavender
%   mygray (near-white) -> lightgray
%   pink                -> Pink
% AAAI comments these out (class conflicts).
\makeatletter
\if@isAAAI\else
  \newcolumntype{a}{>{\columncolor{Lavender}}l}
  % \newcolumntype{b}{>{\columncolor{LimeGreen}}c}
  \DeclareRobustCommand{\hlpurple}[1]{{\sethlcolor{Lavender}\hl{#1}}}
  \DeclareRobustCommand{\hlpink}[1]{{\sethlcolor{Pink}\hl{#1}}}
  \DeclareRobustCommand{\hlgray}[1]{{\sethlcolor{lightgray}\hl{#1}}}
\fi
\makeatother

% \renewcommand{\alertColor}{\textcolor{Red}}
% \renewcommand{\addColor}{\textcolor{Cyan}}

% -------  cleveref  ----------------------------------------
\usepackage[capitalize,noabbrev]{cleveref}
% AAAI comments all crefnames out (conflicts with its own referencing).
\makeatletter
\if@isAAAI\else
  \crefname{ineq}{Inequality}{Inequalities}
  \creflabelformat{ineq}{#2{\upshape(#1)}#3}
\fi
% obj / exp cref names: Elsevier, LNCS, IEEE, ICML, NeurIPS.
\if@isELS
  \crefname{obj}{Objective}{Objectives}
  \creflabelformat{obj}{#2{\upshape(#1)}#3}
  \crefname{exp}{Expression}{Expressions}
  \creflabelformat{exp}{#2{\upshape(#1)}#3}
\fi
\if@isLNCS
  \crefname{obj}{Objective}{Objectives}
  \creflabelformat{obj}{#2{\upshape(#1)}#3}
  \crefname{exp}{Expression}{Expressions}
  \creflabelformat{exp}{#2{\upshape(#1)}#3}
\fi
\if@isIEEE
  \crefname{obj}{Objective}{Objectives}
  \creflabelformat{obj}{#2{\upshape(#1)}#3}
  \crefname{exp}{Expression}{Expressions}
  \creflabelformat{exp}{#2{\upshape(#1)}#3}
\fi
\if@isICML
  \crefname{obj}{Objective}{Objectives}
  \creflabelformat{obj}{#2{\upshape(#1)}#3}
  \crefname{exp}{Expression}{Expressions}
  \creflabelformat{exp}{#2{\upshape(#1)}#3}
\fi
\if@isNIPS
  \crefname{obj}{Objective}{Objectives}
  \creflabelformat{obj}{#2{\upshape(#1)}#3}
  \crefname{exp}{Expression}{Expressions}
  \creflabelformat{exp}{#2{\upshape(#1)}#3}
\fi
\makeatother

% -------  theorem environments  ----------------------------
% reptheorem helper: lets you restate theorems by label.
% Used in: base, ACL, ICLR, USENIX.
% Not used in: ACM, ELS, LNCS, IEEE, ICML, NeurIPS, AAAI.
\makeatletter
\if@isACM\else\if@isELS\else\if@isLNCS\else
\if@isIEEE\else\if@isICML\else\if@isNIPS\else\if@isAAAI\else
  \newtheorem*{rep@theorem}{\rep@title}
  \newcommand{\newreptheorem}[2]{%
  \newenvironment{rep#1}[1]{%
  \def\rep@title{#2 \ref{##1}}%
  \begin{rep@theorem}}%
  {\end{rep@theorem}}}
\fi\fi\fi\fi\fi\fi\fi
\makeatother

% LNCS defines all theorem environments itself — skip them there.
% ICML comments all theorems out — skip them there too.
\makeatletter
\if@isLNCS\else\if@isICML\else

  \theoremstyle{remark}
  \newtheorem{example}{Example}[]
  % reptheorem variants: base, ACL, ICLR, USENIX only
  \if@isACM\else\if@isELS\else\if@isIEEE\else\if@isNIPS\else\if@isAAAI\else
    \newreptheorem{example}{Example}
  \fi\fi\fi\fi\fi

  \theoremstyle{plain}
  \newtheorem{theorem}{Theorem}[section]
  \if@isACM\else\if@isELS\else\if@isIEEE\else\if@isNIPS\else\if@isAAAI\else
    \newreptheorem{theorem}{Theorem}
  \fi\fi\fi\fi\fi

  \newtheorem{proposition}[theorem]{Proposition}
  \newtheorem{lemma}[theorem]{Lemma}
  \if@isACM\else\if@isELS\else\if@isIEEE\else\if@isNIPS\else\if@isAAAI\else
    \newreptheorem{lemma}{Lemma}
  \fi\fi\fi\fi\fi

  \newtheorem{corollary}[theorem]{Corollary}
  \if@isACM\else\if@isELS\else\if@isIEEE\else\if@isNIPS\else\if@isAAAI\else
    \newreptheorem{corollary}{Corollary}
  \fi\fi\fi\fi\fi

  \theoremstyle{definition}
  \newtheorem{definition}[theorem]{Definition}
  % AAAI omits assumption and remark
  \if@isAAAI\else
    \newtheorem{assumption}[theorem]{Assumption}
    \theoremstyle{remark}
    \newtheorem{remark}[theorem]{Remark}
  \fi

\fi\fi % end \if@isLNCS\else\if@isICML\else
\makeatother

% ACM-only: requirement and research question theorem types.
\makeatletter
\if@isACM
  \newtheorem{requirement}{Requirement}[]
  \newtheoremstyle{researchquestionstyle}{}{}{}{}{{\bf}}{}{.5em}{}
  \theoremstyle{researchquestionstyle}
  \newtheorem{rquestion}{Research Question:}[]
\fi
% IEEE-only: requirement (but not rquestion).
\if@isIEEE
  \newtheorem{requirement}{Requirement}[]
\fi
\makeatother

% -------  math operators & notation  -----------------------
\usepackage{braket}
\newcommand{\norm}[1]{\left\lVert#1\right\rVert}
\newcommand{\abs}[1]{\left\lvert#1\right\rvert}
\newcommand{\concat}{\mathbin{{+}\mspace{-8mu}{+}}}
\newcommand{\inner}[1]{\left\langle#1\right\rangle}
\usepackage{stmaryrd}
\newcommand{\iverson}[1]{\left\llbracket#1\right\rrbracket}

% AAAI comments argmax/argmin out (provided by its class).
\makeatletter
\if@isAAAI\else
  \DeclareMathOperator*{\argmax}{arg\,max}
  \DeclareMathOperator*{\argmin}{arg\,min}
\fi
\makeatother

% -------  physics (ACM only)  ------------------------------
\makeatletter
\if@isACM
  \usepackage{physics}
\fi
\makeatother

% -------  TikZ & PGFPlots  ---------------------------------
\usepackage{tikz}
\usetikzlibrary{trees}
\usepackage{tikz-dependency}
\usepackage{tikz-3dplot}
\usetikzlibrary{chains,scopes}
\usetikzlibrary{arrows.meta,automata,positioning}
\usetikzlibrary{shapes.geometric, arrows}
\usetikzlibrary{shapes.geometric, arrows.meta, positioning, calc, fit, shadows}

% ACM-only: enforce LaTeX-style arrow heads globally.
\makeatletter
\if@isACM
  \tikzset{>=latex}
\fi
\makeatother

\usepackage{pgfplots}
% compat=newest + fillbetween: base, ACL, ICLR, USENIX.
% fillbetween only (no compat=newest): AAAI.
% Neither: ACM, LNCS, ELS, IEEE, ICML, NeurIPS.
\makeatletter
\if@isACL
  \pgfplotsset{compat=newest}
  \usepgfplotslibrary{fillbetween}
\else\if@isICLR
  \pgfplotsset{compat=newest}
  \usepgfplotslibrary{fillbetween}
\else\if@isUSENIX
  \pgfplotsset{compat=newest}
  \usepgfplotslibrary{fillbetween}
\else\if@isAAAI
  \usepgfplotslibrary{fillbetween}
\else\if@isACM\else\if@isLNCS\else\if@isELS\else
\if@isIEEE\else\if@isICML\else\if@isNIPS\else
  % base/generic: also use compat=newest + fillbetween
  \pgfplotsset{compat=newest}
  \usepgfplotslibrary{fillbetween}
\fi\fi\fi\fi\fi\fi\fi\fi\fi\fi
\makeatother

\pgfplotsset{
  every axis/.append style = {thick},
  tick style = {thick,black},
  %
  % #1 = x, y, or z  |  #2 = the shift value
  /tikz/normal shift/.code 2 args = {%
    \pgftransformshift{%
        \pgfpointscale{#2}{\pgfplotspointouternormalvectorofticklabelaxis{#1}}%
    }%
  },%
  shift/.style = {
    tick align        = outside,
    scaled ticks      = false,
    enlargelimits     = false,
    ticklabel shift   = {#1},
    axis lines*       = left,
    xtick style       = {normal shift={x}{#1}},
    ytick style       = {normal shift={y}{#1}},
    x axis line style = {normal shift={x}{#1}},
    y axis line style = {normal shift={y}{#1}},
  },
  shift/.default = 10pt,
  shift3d/.style = {
    shift=#1,
    ztick style       = {normal shift={z}{#1}},
    z axis line style = {normal shift={z}{#1}},
  },
  shift3d/.default = 10pt,
}

% TikZ flow-chart styles — ELS, LNCS, IEEE, ICML, NeurIPS, AAAI.
\makeatletter
\if@isELS
  \tikzstyle{startstop} = [rectangle, rounded corners, minimum width=1cm, minimum height=0.5cm, text centered, draw=black]
  \tikzstyle{io}        = [trapezium, trapezium left angle=70, trapezium right angle=110, minimum width=2cm, minimum height=0.63cm, text centered, draw=black]
  \tikzstyle{process}   = [rectangle, minimum width=2cm,   minimum height=0.5cm, text centered, draw=black]
  \tikzstyle{process_n} = [rectangle, minimum width=1.5cm, minimum height=0.5cm, text centered, text width=1.5cm, draw=black]
  \tikzstyle{decision}  = [diamond, minimum width=1.2cm, minimum height=0.5cm, text width=1.5cm, text centered, aspect=1.5, draw=black]
  \tikzstyle{arrow}     = [->,>=stealth]
\fi
\if@isLNCS
  \tikzstyle{startstop} = [rectangle, rounded corners, minimum width=1cm, minimum height=0.5cm, text centered, draw=black]
  \tikzstyle{io}        = [trapezium, trapezium left angle=70, trapezium right angle=110, minimum width=2cm, minimum height=0.63cm, text centered, draw=black]
  \tikzstyle{process}   = [rectangle, minimum width=2cm,   minimum height=0.5cm, text centered, draw=black]
  \tikzstyle{process_n} = [rectangle, minimum width=1.5cm, minimum height=0.5cm, text centered, text width=1.5cm, draw=black]
  \tikzstyle{decision}  = [diamond, minimum width=1.2cm, minimum height=0.5cm, text width=1.5cm, text centered, aspect=1.5, draw=black]
  \tikzstyle{arrow}     = [->,>=stealth]
  % LNCS-only: anchor-highlight helper
  \newcommand{\markline}[1]{%
    \tikz[remember picture,baseline] \node (#1) {};%
  }
\fi
\if@isIEEE
  \tikzstyle{startstop} = [rectangle, rounded corners, minimum width=1cm, minimum height=0.5cm, text centered, draw=black]
  \tikzstyle{io}        = [trapezium, trapezium left angle=70, trapezium right angle=110, minimum width=2cm, minimum height=0.63cm, text centered, draw=black]
  \tikzstyle{process}   = [rectangle, minimum width=2cm,   minimum height=0.5cm, text centered, draw=black]
  \tikzstyle{process_n} = [rectangle, minimum width=1.5cm, minimum height=0.5cm, text centered, text width=1.5cm, draw=black]
  \tikzstyle{decision}  = [diamond, minimum width=1.2cm, minimum height=0.5cm, text width=1.5cm, text centered, aspect=1.5, draw=black]
  \tikzstyle{arrow}     = [->,>=stealth]
\fi
\if@isICML
  \tikzstyle{startstop} = [rectangle, rounded corners, minimum width=1cm, minimum height=0.5cm, text centered, draw=black]
  \tikzstyle{io}        = [trapezium, trapezium left angle=70, trapezium right angle=110, minimum width=2cm, minimum height=0.63cm, text centered, draw=black]
  \tikzstyle{process}   = [rectangle, minimum width=2cm,   minimum height=0.5cm, text centered, draw=black]
  \tikzstyle{process_n} = [rectangle, minimum width=1.5cm, minimum height=0.5cm, text centered, text width=1.5cm, draw=black]
  \tikzstyle{decision}  = [diamond, minimum width=1.2cm, minimum height=0.5cm, text width=1.5cm, text centered, aspect=1.5, draw=black]
  \tikzstyle{arrow}     = [->,>=stealth]
\fi
\if@isNIPS
  \tikzstyle{startstop} = [rectangle, rounded corners, minimum width=1cm, minimum height=0.5cm, text centered, draw=black]
  \tikzstyle{io}        = [trapezium, trapezium left angle=70, trapezium right angle=110, minimum width=2cm, minimum height=0.63cm, text centered, draw=black]
  \tikzstyle{process}   = [rectangle, minimum width=2cm,   minimum height=0.5cm, text centered, draw=black]
  \tikzstyle{process_n} = [rectangle, minimum width=1.5cm, minimum height=0.5cm, text centered, text width=1.5cm, draw=black]
  \tikzstyle{decision}  = [diamond, minimum width=1.2cm, minimum height=0.5cm, text width=1.5cm, text centered, aspect=1.5, draw=black]
  \tikzstyle{arrow}     = [->,>=stealth]
\fi
\if@isAAAI
  \tikzstyle{startstop} = [rectangle, rounded corners, minimum width=1cm, minimum height=0.5cm, text centered, draw=black]
  \tikzstyle{io}        = [trapezium, trapezium left angle=70, trapezium right angle=110, minimum width=2cm, minimum height=0.63cm, text centered, draw=black]
  \tikzstyle{process}   = [rectangle, minimum width=2cm,   minimum height=0.5cm, text centered, draw=black]
  \tikzstyle{process_n} = [rectangle, minimum width=1.5cm, minimum height=0.5cm, text centered, text width=1.5cm, draw=black]
  \tikzstyle{decision}  = [diamond, minimum width=1.2cm, minimum height=0.5cm, text width=1.5cm, text centered, aspect=1.5, draw=black]
  \tikzstyle{arrow}     = [->,>=stealth]
\fi
\makeatother

% -------  IEEE-only: paragraph heading style  --------------
\makeatletter
\if@isIEEE
  \renewcommand\paragraph[1]{{\vspace{1mm}\noindent\bfseries #1 }}
\fi
\makeatother

% -------  font icons (ACM only)  ---------------------------
\makeatletter
\if@isACM
  \usepackage{fontawesome5}
  % \usepackage{emoji}   % requires LuaLaTeX/XeLaTeX
\fi
\makeatother

% -------  code listings  -----------------------------------
% AAAI comments listings out (conflicts with its class).
\makeatletter
\if@isAAAI\else
  \usepackage{listings}
  \lstset{
    language=Python,
  %   numbers=left,
  %   numberstyle=\tiny,
    stepnumber=1,
    numbersep=5pt,
    tabsize=4,
    basicstyle=\scriptsize\ttfamily,
    columns=fullflexible,
    keepspaces,
  %   xleftmargin=\parindent,
    showstringspaces=false
  }
\fi
\makeatother

% JSON language definition (ACM only).
% Color mapping: numb -> Maroon | punct -> BrickRed | delim -> OliveGreen
\makeatletter
\if@isACM
  \lstdefinelanguage{json}{
    basicstyle=\scriptsize\ttfamily,
    numbers=none,
    stepnumber=1,
    numbersep=5pt,
    showstringspaces=false,
    breaklines=true,
    frame=single,
    backgroundcolor=\color{lightgray!30},
    literate=
     *{0}{{{ \color{Maroon}0}}}{1}
      {1}{{{ \color{Maroon}1}}}{1}
      {2}{{{ \color{Maroon}2}}}{1}
      {3}{{{ \color{Maroon}3}}}{1}
      {4}{{{ \color{Maroon}4}}}{1}
      {5}{{{ \color{Maroon}5}}}{1}
      {6}{{{ \color{Maroon}6}}}{1}
      {7}{{{ \color{Maroon}7}}}{1}
      {8}{{{ \color{Maroon}8}}}{1}
      {9}{{{ \color{Maroon}9}}}{1}
      {:}{{{ \color{BrickRed}{:}}}}{1}
      {,}{{{ \color{BrickRed}{,}}}}{1}
      {\{}{{{ \color{OliveGreen}{\{}}}}{1}
      {\}}{{{ \color{OliveGreen}{\}}}}}{1}
      {[}{{{ \color{OliveGreen}{[}}}}{1}
      {]}{{{ \color{OliveGreen}{]}}}}{1},
  }
  % \lstset{language=json}
\fi
\makeatother

% -------  array & table helpers  ---------------------------
% AAAI comments these out.
\makeatletter
\if@isAAAI\else
  \usepackage{array}
  \newcolumntype{H}{>{\setbox0=\hbox\bgroup}c<{\egroup}@{}}
  \newcommand{\RomanNumeralCaps}[1]{\MakeUppercase{\romannumeral #1}}
  \newcommand{\rnumcap}[1]{\MakeUppercase{\romannumeral #1}}
  \newcommand\T{\rule{0pt}{3.4mm}}
  \newcommand\B{\rule[-1.5mm]{0pt}{0pt}}
\fi
\makeatother

% -------  tcolorbox summary environment  -------------------
% ICML and NeurIPS do not define this environment.
\makeatletter
\if@isICML\else\if@isNIPS\else
  \usepackage[most]{tcolorbox}
  \newtcolorbox[auto counter]{summary}[1][]{title={\bfseries Summary~\thetcbcounter},enhanced,
    coltitle=black,
    colback=white,
    top=0.3in,
    attach boxed title to top left=
    {xshift=1.5em,yshift=-\tcboxedtitleheight/2},boxrule=0.5pt, sharp corners,
    fonttitle=\bfseries,
    boxed title style={size=small,colback=white,colframe=white},#1}
\fi\fi
\makeatother
